\subsubsection[Recursive arm-specific prior paths]{\color{revisionred}Recursive arm-specific prior paths}
\label{subsec:arm-specific-prior-dose}
{\color{revisionred}
Figure~\ref{fig:arm-specific-prior-dose-gradient} and Table~\ref{tab:arm-specific-prior-dose-arms} report the complementary recursive-path exercise. Unlike the common-prior design, only the initial 3.80 percent forecast is shared. From the second meeting onward, each textual arm receives the forecast generated by that same arm in the preceding meeting; the original text follows its own recursive reference path. The comparison therefore captures accumulated synthetic state dependence rather than the one-shot effect of the current communication.
\par}
\begin{figure}[htbp]
\color{revisionred}
\centering
\includegraphics[width=\textwidth]{graficos/rr_arm_specific_prior_dose_gradient.pdf}
\caption{Generated-text tone and persistent synthetic forecast paths under arm-specific recursive priors}
\label{fig:arm-specific-prior-dose-gradient}
\begin{minipage}{0.95\textwidth}
\footnotesize
\textit{Notes:} Panel A reports the tone score measured in each LLM-generated rewrite, not the tone of the published Copom minute. Panel B reports the calibrated synthetic forecast-level gap relative to the recursive original-text path. Every arm starts from 3.80 percent at the first meeting and then uses only its own lagged forecast as its prior. One basis point equals 0.01 percentage point. Whiskers are 95 percent confidence intervals based on two-way clustered standard errors by Copom meeting and recursive trajectory (regeneration replicate by textual arm).
\end{minipage}
\end{figure}
{\color{revisionred}
Relative to the recursive original-text path, the composite forecast gaps are $+2.78$ bp under high hawk, $+4.17$ bp under low hawk, $-4.29$ bp under neutral, $-4.45$ bp under low dove, and $-4.24$ bp under high dove. 
The linear path gradient is $+2.27$ bp per step toward hawkishness (two-way clustered standard error 0.33; $p<0.001$). The five composite means are not monotone at every adjacent step, so the linear coefficient is interpreted as an average path gradient rather than as evidence that every neighboring intensity is separately ordered. Table~\ref{tab:arm-specific-prior-dose-arms} reports the anchored and unanchored components using the same specification.
\par}
\input{tabelas/rr_arm_specific_prior_dose_arm_effects.tex}
{\color{revisionred}
Because the prior is endogenous to earlier synthetic responses within each path, these coefficients are persistent model-trajectory comparisons. They are deliberately reported alongside, and not substituted for, the immediate common-prior Prompt-IV estimand.
\par}
\FloatBarrier
